\section{RQ1: Do current adaptation methods generalise?}
\label{sec:rq1}

Adapting a code model to obfuscated input works, and works easily: tuning on output prediction over
obfuscated programs gains roughly \textbf{19 points} over the untuned model on the transforms it was
trained on. We apply the two tests of Section~\ref{sec:setup} to the field's existing answers.
These are baselines, not proposals. Merging and breadth fail both tests; routing tracks breadth
wherever the two are measured together and inherits its verdict.

Routing composes exactly as far as breadth does, and merging does not compose at all
(Table~\ref{tab:rq1_composition}, Appendix~\ref{app:supp}). On the ten stacks of seen transforms, a learned router beats a
\emph{random} gate by $+1.34$ [$+0.48$, $+2.17$]: the mixture is worth the whole gain ($+19.79$ over
the untuned model) and the routing about a point of it. A hard argmax router is indistinguishable
from the soft one ($+0.03$ [$-0.13$, $+0.19$]), so a stacked input gives the mixture no blending
advantage, and the router matches breadth training, $-0.28$ [$-1.57$, $+1.09$]. The best merge
(DARE-TIES) matches the clean-code control, $+0.52$ [$-0.15$, $+1.20$], the others fall $1.7$ to
$20.6$ points below it, while breadth stands $+3.67$ [$+2.02$, $+5.24$] above the same control.


% ---------------------------------------------------------------------------------------------
% FRAMING DECISION PENDING (2026-09-13, the user's call). THIRD contradiction: the abstract and this
% section say every method that fits the forward task loses backward competence; merge_dare_ties matches
% breadth forward (+18.87 vs +18.68) and GAINS +3.21 [+1.52, +4.85] backwards, direction ratio +0.17 --
% the highest of any arm (H-merge-backward REFUTED, log/modularity/2026-09-13_merging-gains-backward.md).
% Also: the text says routing "tracks breadth
% wherever the two are measured together"; the fourth column of Table~\ref{tab:rq1_composition} shows the
% mixture +2.88 [+1.33, +4.41] ABOVE breadth on stacks containing the unseen family (H-F2-route REFUTED,
% log/modularity/2026-09-13_f2-route-read.md). The text above also says merging "does not
% compose at all" while Table~\ref{tab:rq1_merge_panel} shows the best merge above the clean-code
% control on seen stacks on 3/3 models and above breadth on unseen-containing stacks on 3/3. Either
% the merging claim is narrowed to the measurement, or the table moves to the appendix; see
% log/modularity/2026-09-13_merging-on-the-panel.md.
% ---------------------------------------------------------------------------------------------
Table~\ref{tab:rq1_merge_panel} (Appendix~\ref{app:supp}) reports the same merge arms measured across the panel.


Breadth is the interesting case. It composes on stacks built from transforms it saw, and the
advantage evaporates the moment a stack contains a family it did not. Table~\ref{tab:rq1_dissociation}
states this as a \emph{difference}, because one interval being positive while another is negative
is not a test of the difference between them: breadth's advantage over clean-code tuning is
\textbf{2.7 to 5.3 points larger} on seen stacks than on stacks containing an unseen family, and the
gap is significant on \textbf{every one of the six models} we could measure it on, across four
lineages. What breadth buys is recombination of what the training data covered, bounded by that
coverage.

% >>> GENERATED TABLE rq1_dissociation -- from tables/rq1_dissociation.tex; regenerate, then re-run inline_tables.py
% GENERATED 2026-09-14 02:18 UTC by scripts/paper/gen_fse_tables.py
% SOURCE: results/analysis/pipeline/stack_dissociation.json
% Do not edit by hand -- regenerate. Supersedes the hand-written '+3.49 / +2.90 / +3.05 across three model scales'.
\begin{table*}[ht]
\centering\footnotesize\setlength{\tabcolsep}{3pt}
\caption{\textbf{RQ1: the dissociation, measured within model on one program set.} Breadth
(\texttt{mono\_all}) against the clean-code control (\texttt{tuned\_L0}) on stacks built from
\emph{seen} transforms and on stacks containing an \emph{unseen} family. The \emph{difference} is
the measured quantity; the two levels are context, and on most models the unseen-side interval
straddles zero on its own. Both halves are paired on the 319 programs common
to every cell of both groups. Bold intervals exclude zero.}
\label{tab:rq1_dissociation}
\begin{tabular}{@{}llccc@{}}
\toprule
\textbf{Model} & \textbf{Lineage} & \textbf{Seen stacks} & \textbf{Unseen inside} & \textbf{Difference} \\
\midrule
CodeLlama-7B & Meta & \textbf{$+3.36$ [$+1.57$, $+5.27$]} & \textbf{$-1.92$ [$-3.70$, $-0.14$]} & \textbf{$+5.28$ [$+3.20$, $+7.54$]} \\
CodeLlama-34B & Meta & \textbf{$+2.27$ [$+0.44$, $+4.09$]} & $-0.77$ [$-2.62$, $+1.12$] & \textbf{$+3.03$ [$+0.78$, $+5.28$]} \\
Llama-3.1-8B & Meta & \textbf{$+2.30$ [$+0.47$, $+4.15$]} & $-1.05$ [$-2.72$, $+0.63$] & \textbf{$+3.35$ [$+1.50$, $+5.30$]} \\
StarCoder2-15B & BigCode & $+1.50$ [$-0.02$, $+3.00$] & $-1.19$ [$-2.93$, $+0.49$] & \textbf{$+2.68$ [$+0.77$, $+4.65$]} \\
Gemma-3-12B & Google & \textbf{$+2.63$ [$+0.56$, $+4.64$]} & \textbf{$-2.09$ [$-3.81$, $-0.42$]} & \textbf{$+4.72$ [$+2.62$, $+6.85$]} \\
Granite-3.1-8B & IBM & \textbf{$+3.66$ [$+1.54$, $+5.68$]} & $-0.24$ [$-2.02$, $+1.57$] & \textbf{$+3.91$ [$+1.41$, $+6.32$]} \\
\midrule
\multicolumn{4}{@{}l}{\emph{difference significant on}} & 6 of 6 models \\
\bottomrule
\end{tabular}
\end{table*}
% <<< GENERATED TABLE rq1_dissociation

The reversal test is independent and gives the same answer. Table~\ref{tab:direction_ratio} (Appendix~\ref{app:supp}) reports
the \emph{direction ratio} --- backward gain divided by forward gain, both against the untuned model
--- for the same adapters on the same programs asked backwards. \textbf{Every method that fits the
forward task directly is negative}: the format-only control at $-0.30$, clean-code tuning at
$-0.08$, breadth at $-0.05$; these arms bought forward accuracy by losing backward competence on the
very same programs. Both tests point at one explanation: what is acquired is competence at \emph{this task on these
surfaces}. Three further instruments, sharing no machinery, say which surface. Breadth's advantage on
a stack depends on the stack still containing an \emph{identifier} transform rather than on its depth
or the number of mechanisms in it, and the rule defining that split was fixed before the
unseen-containing stimulus existed, which makes its second row an out-of-sample test: the difference
is $+3.39$ on all-seen stacks and grows to $+5.85$ on stacks containing the unseen family. Splitting
the held-out family into its two mechanisms, the halves \textbf{do not transfer to each other}
($+1.49$ [$-0.95$, $+3.93$] and $-0.10$ [$-2.38$, $+2.19$]) yet \textbf{either alone recovers about
90\,\%} of the whole adapter's gain on the full family --- composable skills would deliver half each.
And breadth's log-probability of the gold answer sits about 4 nats below the clean-code control with
the entire deficit in the items it gets wrong: it does not become a worse reader, it becomes a more
committed guesser. Adaptation anchors on the surface a transform leaves behind, which is exactly what
an obfuscator is free to change. The remedy follows: anchor the model to something the transform
\emph{cannot} change --- its own behaviour on the unobfuscated parent, which exists for every
training row by construction. Section~\ref{sec:loss} is that objective.

% >>> GENERATED TABLE rq2_identifier -- from tables/rq2_identifier.tex; regenerate, then re-run inline_tables.py
% GENERATED 2026-09-14 02:18 UTC by scripts/paper/gen_fse_tables.py
% SOURCE: results/analysis/pipeline/f3a_stack_identifier_{depth_seen,f2_unseen}_codellama7b.json
% Do not edit by hand -- regenerate. Supersedes the hand-written '+6.6 down to +1.1', which was an ordering, not a contrast.
\begin{table*}[ht]
\centering\footnotesize\setlength{\tabcolsep}{3pt}
\caption{\textbf{RQ2: breadth's stack gain depends on an identifier transform being present.}
Breadth over the clean-code control, split by whether the stack contains an identifier transform,
with the \emph{difference} between the groups. The rule defining the split was fixed before the
unseen-containing stimulus existed, which makes the second row an out-of-sample test of it. We
report the difference because one interval excluding zero while another does not is not a test of
the difference between them.}
\label{tab:rq2_identifier}
\begin{tabular}{@{}lccc@{}}
\toprule
\textbf{Stimulus} & \textbf{With identifier} & \textbf{Structural only} & \textbf{Difference} \\
\midrule
All-seen depth-3/4 stacks (in-sample) & \textbf{$+5.00$ [$+3.02$, $+6.88$]} & $+1.61$ [$-1.02$, $+4.23$] & \textbf{$+3.39$ [$+0.82$, $+6.16$]} \\
Stacks containing the unseen family (out-of-sample) & $+1.98$ [$-0.12$, $+4.18$] & \textbf{$-3.88$ [$-6.08$, $-1.71$]} & \textbf{$+5.85$ [$+3.10$, $+8.72$]} \\
\bottomrule
\end{tabular}
\end{table*}
% <<< GENERATED TABLE rq2_identifier

\begin{takeaway}{Takeaway: RQ1}
Routing, merging and breadth are baselines, and none survives leaving its training distribution
intact.
Breadth's advantage over clean-code tuning is 2.7 to 5.3 points larger on stacks of seen transforms
than on stacks containing an unseen family, significant on six of six models across four lineages;
and every method that fits the forward task has a negative direction ratio, buying forward accuracy
with backward competence on the same programs. Three independent instruments locate what is learned
instead: the surface a transform leaves behind --- which is exactly what an obfuscator is free to
change, and which names the remedy.
\end{takeaway}
